%! The Victory of Orthogonality — Unit 7, Lesson 7.4 — Student Packet Cover Sheet
\documentclass[10pt]{article}
\usepackage{linalg-article}
\usepackage{linalg-boxes}
\usepackage{ltablex}
\keepXColumns

\begin{document}

% Full-bleed burgundy banner
\begin{tikzpicture}[remember picture, overlay]
  \fill[burgundy] (current page.north west)
    rectangle ([yshift=-0.9in]current page.north east);
\end{tikzpicture}
\vspace{-0.8in}

\begin{center}
  {\color{white}\textbf{\LARGE Linear Algebra}}\\[4pt]
  {\color{blushmid}\large\textbf{Unit 7 \quad The Singular Value Decomposition}}\\[6pt]
  {\color{white}\textbf{\large Lesson 7.4 \quad The Victory of Orthogonality}}
\end{center}

\vspace{0.22in}
\namedateperiod
\vspace{0.18in}

\begin{learningtargetbox}
\small
\begin{enumerate}[leftmargin=1.5em, itemsep=5pt]
  \item \ldots{} recognize an \textbf{orthogonal matrix} $Q$ by the single test $Q^{\top}Q=I$ (perpendicular unit columns), and read off its free inverse $Q^{-1}=Q^{\top}$.
  \item \ldots{} show that an orthogonal matrix \textbf{preserves length} ($|Q\mathbf{x}|=|\mathbf{x}|$) --- it moves vectors rigidly, without stretching.
  \item \ldots{} read the SVD $A=U\Sigma V^{\top}$ as \textbf{rotate--stretch--rotate} and explain why orthogonal matrices are the idea behind every result in this unit.
\end{enumerate}
\end{learningtargetbox}

\vspace{0.14in}

\begin{tocbox}
\small
\renewcommand{\arraystretch}{1.5}
\begin{tabularx}{\linewidth}{@{} c l X r @{}}
\rowcolor{blush}
\textbf{\#} & \textbf{Component} & \textbf{Description} & \textbf{Score} \\
\midrule
1 & Warm-Up        & Spiral review: perpendicular unit columns, $Q^{\top}Q=I$, and length after $Q$      & \blank{1.2cm} \\
2 & Guided Notes   & Orthogonal matrices, the free inverse, length preservation, rotate--stretch--rotate & \blank{1.2cm} \\
3 & Group Activity & Verify a matrix is orthogonal, show length is preserved, read the SVD's factors     & \blank{1.2cm} \\
4 & Exit Ticket    & Quick check: verify orthogonality, the free inverse, and justify                    & \blank{1.2cm} \\
5 & Homework       & Orthogonal-or-not, length \& free inverse, the SVD factors, and a $\det=\pm1$ fact   & \blank{1.2cm} \\
\midrule
  & \multicolumn{2}{r}{\textbf{Total}} & \blank{1.2cm} \\
\end{tabularx}
\end{tocbox}

\vspace{0.10in}

\begin{remindbox}
\small An \textbf{orthogonal matrix} $Q$ is a square matrix whose columns are \textbf{perpendicular unit
vectors} --- captured by the one equation $Q^{\top}Q=I$. Its inverse is then \emph{free}: $Q^{-1}=Q^{\top}$,
and it \textbf{preserves length}, $|Q\mathbf{x}|=|\mathbf{x}|$, so it moves space rigidly (a rotation or a
reflection). The SVD $A=U\Sigma V^{\top}$ writes \emph{every} matrix using two orthogonal matrices and a
diagonal one: $V^{\top}$ turns, $\Sigma$ stretches by the singular values, $U$ turns --- \textbf{rotate,
stretch, rotate}. This is orthogonality's victory: the reason the SVD, image compression, and PCA all work.
\end{remindbox}

\end{document}
